\documentclass[a4paper,10pt]{report}\input{../0.Préambule/Préambule_cours_prof}\begin{document}

\newpage
\section*{Statistiques}
\addcontentsline{toc}{section}{Statistiques}

\begin{exer1}
On considère deux séries de notes.
\begin{enumerate}[font=\color{exer1}]
\begin{tabularx}{\linewidth}{*{2}{>{\arraybackslash}X}}
\item[•]Série A : $12$ ; $8$ ; $16$ ; $10$.&
\item[•]Série B : $8$ ; $13$ ; $11$ ; $9$ ; $14$.\\
\end{tabularx}
\end{enumerate}
Calculer la moyenne des notes de chaque série.
\end{exer1}

\begin{exer1}
Dans une entreprise, le patron se vante de respecter l'égalité hommes-femmes. Il dit employer $6$ hommes et $6$ femmes. Cependant, l'égalité ne s'arrête pas à la quantité mais aux salaires, on a récupéré les salaires de chaque personnel de l'entreprise. Les salaires sont rangés dans l'ordre croissant et par genre. Voici le tableau récapitulatif.
\begin{center}
\renewcommand{\arraystretch}{1.5}
\begin{tabularx}{\linewidth}{|>{\columncolor{exer1!20}\centering}p{2cm}|*{6}{>{\centering \arraybackslash}X|}}
\hline
Hommes & $1 725$\euro & $1 943$\euro & $1 945$\euro & $2 027$\euro & $3 121$\euro & $4 235$\euro \\\hline
Femmes & $1 324$\euro & $2 002$\euro & $2 012$\euro & $2 035$\euro & $2 204$\euro & $3 126$\euro \\\hline
\end{tabularx}
\end{center}	

\begin{enumerate}[font=\color{exer1}]
\item Déterminer le salaire moyen chez les hommes puis chez les femmes.
\item Code du travail : \og Tout employeur est tenu d'assurer, pour un même travail ou un travail de valeur égale, l'égalité de rémunération entre les femmes et les hommes : ce principe interdit toute discrimination de salaire fondée sur le sexe. Tous les employeurs et tous les salariés sont concernés, qu'ils relèvent ou non du Code du travail. Les salariés du secteur public sont donc également visés. \fg{}

\medskip
\noindent Est-ce que le patron respecte cette loi ?
\item Pour attirer de nouveaux employé$\cdot$e$\cdot$s, le patron dit : \og Le salaire moyen de mon entreprise est de $2 308$\euro donc $50 \%$ de mes employé$\cdot$e$\cdot$s touchent plus que ce salaire \fg{}. Qu'en pensez vous ?
\end{enumerate} 
\end{exer1}

\begin{exer1}
Voici le tableau représentant l'espérance de vie en France à la naissance selon le genre.

\medskip
\noindent \begin{minipage}{8cm}
\begin{center}
\renewcommand{\arraystretch}{1.3}
\begin{tabularx}{\linewidth}{|*{3}{>{\centering \arraybackslash}X|}}
\hline
\rowcolor{exer1!20}Année & Espérance de vie des femmes & Espérance de vie des hommes\\\hline
$2020$ & $85,2$ & $79,2$ \\\hline
$2019$ & $85,6$ & $79,7$ \\\hline
$2018$ & $85,4$ & $79,5$ \\\hline
$2017$ & $85,3$ & $79,4$ \\\hline
$2010$ & $84,6$ & $78,0$ \\\hline
\end{tabularx}
\end{center}
\end{minipage}
\hspace{5mm}
\begin{minipage}{8cm}
\textbf{Définition de l'espérance de vie :} L'espérance de vie à la naissance représente la durée de vie moyenne d'une génération fictive soumise aux conditions de mortalité par âge de l'année considérée.

\vspace{-5mm}
\begin{flushright}
\textit{Source : INSEE}
\end{flushright}

\begin{enumerate}[font=\color{exer1}]
\item Que peut-on dire de l'espérance de vie en France entre les hommes et les femmes ?
\item Déterminer l'espérance de vie moyenne chez les femmes entre 2010 et 2020.
\item Faire de même pour les hommes.
\end{enumerate}
\end{minipage}
\end{exer1}

\begin{exer2}
On s'intéresse à l'empreinte carbone (E.C.) par personne sur le territoire national en France. Voici les données selon les années :

\begin{center}
\renewcommand{\arraystretch}{1.2}
\begin{tabularx}{\linewidth}{|>{\columncolor{exer2!20}}p{3cm}|*{8}{>{\centering \arraybackslash}X|}}
\hline
Années & $2012$ & $2013$ & $2014$ & $2015$ & $2016$ & $2017$ & $2018$ & $2019$\\\hline
E.C (en  tonnes) & $11,4$ & $11,0$ & $10,8$ & $10,5$ & $10,0$ & $10,1$ & $9,7$ & $9,9$\\\hline
\end{tabularx}
\end{center}

\begin{enumerate}[font=\color{exer2}]
\item Calculer les émissions moyenne de carbone pour un habitant en France entre $2012$ et $2019$.

\item L'accord de Paris signé en $2016$ demande à l'état Français de réduire significativement ses émissions de CO$_2$. Le rapport est-il respecté ?
\end{enumerate}
\end{exer2}

\begin{exer2}
L'état Français finance actuellement l'éducation des enfants sur le territoire français. On s'intéresse à la dépense moyenne pour un élève du second degré (collège + lycée) selon les années.

\begin{center}
\renewcommand{\arraystretch}{1.2}
\begin{tabularx}{\linewidth}{|>{\columncolor{exer2!20}}p{3cm}|*{7}{>{\centering \arraybackslash}X|}}
\hline
Année & $2019$ & $2018$ & $2017$ & $2016$ & $2015$ \\\hline
Dépense en euros & $9 950$ & $10 000$ & $10 040$ & $9 940$ & $9 930$\\\hline
\end{tabularx}
\end{center}

\begin{enumerate}[font=\color{exer2}]
\item Calculer la moyenne des dépenses pour un élève du second degré entre 2019 et 2015.
\item Que peut-on dire de l'évolution de ses dépenses ?
\end{enumerate}
\end{exer2}

\begin{exer1}
Calcule, à l'aide de ta calculatrice, la moyenne arrondie au dixième de la série :

\begin{center}
\renewcommand{\arraystretch}{1.2}
\begin{tabularx}{\linewidth}{|>{\columncolor{exer1!20}}p{2cm}|*{7}{>{\centering \arraybackslash}X|}}
\hline
Valeurs	& $26$ & $33$ & $152$ & $45$ & $89$ & $78$ & $45$\\\hline
Coefficient & $2$ & $5$ & $3$ & $4$ & $8$ & $10$ & $6$\\\hline
\end{tabularx}
\end{center}
\end{exer1}

\begin{exer1}
Dans une classe, on relève la durée, en minutes, du trajet maison-collège. Les données, par élève, sont les suivantes :

\noindent $30$ \quad $45$ \quad $10$ \quad $30$ \quad $50$ \quad $20$ \quad $25$ \quad $25$ \quad $60$ \quad $30$ \quad $20$ \quad $25$ \quad $20$ \quad $25$ \quad $5$ \quad $10$ \quad $45$ \quad $30$ \quad $20$ \quad $25$ \quad $5$ \quad $10$ \quad $25$ \quad $45$

\begin{enumerate}[font=\color{exer1}]
\item Complète le tableau statistique suivant (les valeurs de la série seront rangées dans l'ordre croissant) :

\begin{center}
\renewcommand{\arraystretch}{1.5}
\begin{tabularx}{\linewidth}{|>{\columncolor{exer1!20}}p{2cm}|*{8}{>{\centering \arraybackslash}X|}}
\hline
Durée & & & & & & & &\\\hline							
Effectif & & & & & & & &\\\hline
\end{tabularx}
\end{center}
						
\item Calcule la durée moyenne du trajet des élèves de cette classe.
\end{enumerate}
\end{exer1}

\begin{exer1}
Dans un groupe de personnes, on considère le nombre de frères et sœurs. On relève les données statistiques dans le tableau suivant :

\begin{center}
\renewcommand{\arraystretch}{1.2}
\begin{tabularx}{\linewidth}{|>{\columncolor{exer1!20}}p{4cm}|*{8}{>{\centering \arraybackslash}X|}}
\hline
Nombre de frères et sœurs & $0$ & $1$ & $2$ & $3$ & $4$ & $5$ & $6$ & $7$\\\hline
Effectif & $3$ & $6$ & $7$ & $9$ & $5$ & $2$ & $1$ & $1$\\\hline
\end{tabularx}
\end{center}

\begin{enumerate}[font=\color{exer1}]
\item Donne l'effectif total de cette série.
\item Combien de personnes ont quatre frères et sœurs ? Combien de personnes ont au moins trois frères et sœurs ?
\item Calcule le nombre moyen de frères et sœurs.
\end{enumerate}
\end{exer1}

\begin{exer2}
En Mathématiques, Adélaïde a des notes de contrôles en classe (coefficient $7$) et des notes de devoirs maison (coefficient $1$).
Voici les notes d'Adélaïde pour un trimestre :

\begin{tabularx}{\linewidth}{*{2}{>{\arraybackslash}X}}
$\bullet$ En contrôle : $7$ \quad $9$ \quad $11$ \quad $9,5$ \quad $10,5$ \quad $8$ &

$\bullet$ En devoir maison : $13$ \quad $14$ \quad $12$ \quad $11$ \\
\end{tabularx}

\begin{enumerate}[font=\color{exer2}]
\item Recopie et complète la phrase suivante :

\medskip
\noindent \og La note $7$ en contrôle compte $\dots ~\dots ~\dots$ fois dans la moyenne.\fg{}
\item Pour calculer sa moyenne du trimestre, par quel nombre faut-il diviser ? Calculer cette moyenne.
\item Pour augmenter sa moyenne, est-il préférable d'avoir $3$ points de plus à un devoir maison ou $2$ points de plus à un contrôle ?
\end{enumerate}
\end{exer2}

\begin{exer2}
L'âge et la taille des joueurs d'une équipe de football (avec les remplaçants) sont précisés ci-dessous.

\begin{center}
\renewcommand{\arraystretch}{1.2}
\begin{tabularx}{\linewidth}{|>{\columncolor{exer2!20}}p{2cm}|*{9}{>{\centering \arraybackslash}X|}}
\hline
Age & $33$ & $27$ & $27$ & $23$ & $33$ & $27$ & $20$ & $27$ & $33$	\\\hline
Taille (en m) & $1,89$ & $1,74$ & $1,87$ & $1,86$ & $1,70$ & $1,72$ & $1,79$ & $1,76$ & $1,78$\\\hline\hline
Age & $27$ & $27$ & $20$ & $27$ & $23$ & $23$ & $21$ & $23$ & $27$\\\hline
Taille (en m) & $1,75$ & $1,80$ & $1,83$ & $1,76$ & $1,74$ & $1,83$ & $1,71$ & $1,88$ & $1,72$\\\hline
\end{tabularx}
\end{center}

\begin{enumerate}[font=\color{exer2}]
\item Calculer l'âge moyen des joueurs.
\item \begin{enumerate}[font=\color{exer2}]
\item Réaliser le tableau des effectifs des âges des joueurs.
\item Utiliser ce tableau pour calculer d'une autre façon l'âge moyen des joueurs.
\end{enumerate}
\item Représenter la répartition des joueurs selon leur âge à l'aide d'un diagramme en bâtons.
\item \begin{enumerate}[font=\color{exer2}]
\item Calculer la taille moyenne des joueurs.
\item Est-il avantageux de construire un tableau des effectifs pour calculer cette taille ?

\noindent Expliquer pourquoi.
\end{enumerate}
\end{enumerate}
\end{exer2}

\begin{exer3}
Voici les notes obtenues par Aurélie pendant une année en mathématiques :

\begin{center}
\renewcommand{\arraystretch}{1.2}
\begin{tabularx}{\linewidth}{|>{\columncolor{exer3!20}}p{2cm}|*{7}{>{\centering \arraybackslash}X|}}
\hline
T1 & $10$ & $9$ & $11$ & $12$ & $11,5$ & $14$ & $12$ \\\hline
T2 & $9,5$ & $11$ & $12,5$ & $8$ & $13$ & $14$ &  \\\hline
T3 & $7,5$ & $9$ & $14$ & $12$ & $10$ & $13$ & $11,5$ \\\hline
\end{tabularx}
\end{center}
Toutes les notes ont le même coefficient.

\begin{enumerate}[font=\color{exer3}]
\item Calcule la moyenne de toutes les notes de l'année.
\item Calcule la moyenne de chaque trimestre.
\item Calculer la moyenne des moyennes trimestrielles. Compare-la avec la première moyenne calculée. Que peut-on dire de ces deux résultats ? Pourquoi ?
\item Construis une série de note de manière à ce que la moyenne des notes de l'année soit supérieure à la moyenne des trois trimestres.
\item Construis une série de notes de manières à ce que la moyenne des notes de l'année soit inférieur à la moyenne des trois trimestres.
\end{enumerate}
\end{exer3}

\begin{exer1}
Ce tableau fournit les températures mensuelles moyennes (en °C) au cours d’une année dans deux villes Alpha (A) et Gamma (G).

\begin{center}
\renewcommand{\arraystretch}{1.2}
\begin{tabularx}{\linewidth}{|>{\columncolor{exer1!20}}p{0.5cm}|*{12}{>{\centering \arraybackslash}X|}}
\hline
\rowcolor{exer1!20} & J & F & M & A & M & J & J & A & S & O & N & D \\\hline
A & $-6$ & $-9$ & $-1$ & $10$ & $11$ & $19$ & $24$ & $28$ & $21$ & $10$ & $4$ & $-3$ \\\hline
G & $5$ & $7$ & $9$ & $13$ & $17$ & $19$ & $20$ & $23$ & $18$ & $13$ & $8$ & $4$ \\\hline
\end{tabularx}
\end{center}

Pour la ville Alpha puis pour la ville Gamma :
\begin{enumerate}[font=\color{exer1}]
\item Calcule la moyenne des températures.
\item Détermine une médiane des températures.
\item Calcule l'étendue des températures.
\end{enumerate}
\end{exer1}

\begin{exer1}
Un gérant a relevé le nombre de personnes fréquentant son club sur une semaine.
\begin{center}
\renewcommand{\arraystretch}{1.2}
\begin{tabularx}{\linewidth}{|*{7}{>{\centering \arraybackslash}X|}}
\hline
\rowcolor{exer1!20} Lu & Ma & Me & Je & Ve & Sa & Di \\\hline
$32$ & $38$ & $21$ & $49$ & $60$ & $84$ & $24$\\\hline
\end{tabularx}
\end{center}

\begin{enumerate}[font=\color{exer1}]
\item Calcule le nombre moyen de personnes fréquentant le club par jour.
\item Détermine une médiane de cette série.
\end{enumerate}
\end{exer1}

\begin{exer1}
Détermine la médiane de chacune des séries statistiques suivantes.

\begin{enumerate}[font=\color{exer1}]
\item $10$ - $25$ – $32$ – $143$ – $280$ – $301$
\item $8$ – $9.5$ – $14$ – $16$ – $18$ – $19$ – $20$
\item $2$ – $4$ – $7$ – $9$ – $9$ – $10$ – $10$ – $10$ – $11$ – $12$
\item $1$ – $1$ – $1$ – $3$ – $4$ – $4$ – $5$ – $5$ – $5$ – $9$ – $10$
\end{enumerate}
\end{exer1}

\begin{exer2}
Ce tableau donne la répartition des salaires mensuels des employés d’une petite entreprise.

\begin{center}
\renewcommand{\arraystretch}{1.2}
\begin{tabularx}{\linewidth}{|>{\columncolor{exer2!20}}p{3cm}|*{5}{>{\centering \arraybackslash}X|}}
\hline
Salaire(en \euro) & $1 000$ à $1 200$ & $1 200$ à $1 400$ & $1 400$ à $1 600$ & $1 600$ à $1 800$ & $2 000$ à $2 200$ \\\hline
Fréquence(en $\%$) & $6,5$ & $9,5$ & $38,5$ & $25,5$ & $20$ \\\hline
\end{tabularx}
\end{center}

\begin{enumerate}[font=\color{exer2}]
\item Calcule une valeur approchée du salaire moyen d'un employé.
\item Dans quelle classe est situé le salaire médian ? Que signifie-t-il ?
\end{enumerate}
\end{exer2}

\begin{exer3}
Sam a relevé les durées des morceaux de sa compilation de rap préférée en min:sec.

\medskip
\noindent $4$:$08$ \quad ; \quad $3$:$19$ \quad ; \quad $4$:$47$ \quad ; \quad $3$:$46$ \quad ; \quad $3$:$15$ \quad ; \quad $3$:$19$ \quad ; \quad $3$:$58$ \quad ; \quad $3$:$50$ \quad ; \quad $3:24$ \quad ; \quad $3$:$55$ \quad ; \quad $3$:$16$ \quad ; \quad $3$:$24$ \quad ; \quad $3$:$07$ 

\begin{enumerate}[font=\color{exer3}]
\item Calcule la durée moyenne des morceaux.
\item Détermine la durée médiane.
\end{enumerate}
\end{exer3}

\begin{exer3}

\medskip
\begin{enumerate}[font=\color{exer3}]
\item Donner l'exemple d'une série statistique où la moyenne est négative.
\item Donner l'exemple d'une série statistique où la moyenne et la médiane est égale à $0$.
\end{enumerate}
\end{exer3}


\end{document}